\chapter{Visão computacional}\label{cap:visao}
A visão computacional compõe a observação associada à tomada, em complemento aos registros do mecanismo e aos relatos do cuidado. Sua função é produzir uma saída rastreável sobre as imagens analisadas, preservando a condição inconclusiva quando não houver informação suficiente. Uma observação visual não é tratada como comprovação clínica de ingestão.

A integração com a central compartilha o recurso de câmera e precisa delimitar aquisição, processamento e devolução do resultado. O registro deve conservar o evento de referência, a versão do método e as condições da captura, de forma que a saída possa ser relacionada à rotina apresentada no aplicativo.

\section{Método de aquisição e processamento}
\pendencia{VIS01}{Completar a descrição técnica da frente de visão com o objetivo operacional e o evento observado. Responsável: Vinicius.}
\pendencia{VIS02}{Inserir a arquitetura implementada: aquisição, preparação, algoritmo ou modelo, saídas, erros, versões e implantação. Acrescentar o diagrama e as imagens do método utilizado.}
\section{Dados e avaliação}
\pendencia{VIS03}{Informar origem e autorização dos dados, anonimização, classes, rotulagem, composição e separação dos conjuntos e referência independente.}
\pendencia{VIS04}{Definir métricas, limiares, cenários, repetições, condições e critérios de aceitação. Distinguir a pontuação do modelo de certeza clínica.}
\section{Resultados e limitações}
\pendencia{VIS05}{Inserir somente os resultados experimentais efetivamente obtidos, com versões, denominadores e registros correspondentes.}
\pendencia{VIS06}{Descrever as falhas observadas, fatores de generalização e condições em que o método deve declarar resultado inconclusivo.}
